% Осень26!A31; нумерация внутри столбца без пустых строк.
% Источник: лекция 6, §1 (строка 26)
% Источник: лекция 6, §2 (строка 52)
\ticket{28}{Простейшие свойства интеграла Лебега}
\label{ticket:28}

\begin{definition}
Для измеримого $E\subset X$ его \textbf{индикатор} определяется формулой
\[
  \chi_E(x)=\begin{cases}1,&x\in E,\\0,&x\notin E.\end{cases}
\]
\textbf{Интеграл по множеству} задаётся равенством
\[
  \int_E f\,d\mu:=\int_X f\chi_E\,d\mu.
\]
Если ровно один из интегралов $\int f^+\,d\mu$, $\int f^-\,d\mu$ бесконечен, формулу из \hyperref[definition:26]{О26} также используют: она определяет интеграл со значением $+\infty$ или $-\infty$.
\end{definition}

\begin{theorem}
Пусть $(X,\mathcal A,\mu)$ --- пространство с мерой, $E\in\mathcal A$; все функции ниже измеримы, а суммируемость рассматривается на $E$.
\begin{enumerate}
  \item Если $f\le g$ и обе функции суммируемы либо обе неотрицательны, то $\int_E f\,d\mu\le\int_E g\,d\mu$.
  \item $\int_E1\,d\mu=\mu(E)$ и $\int_E0\,d\mu=0$.
  \item Если $\mu(E)=0$, то $\int_E f\,d\mu=0$.
  \item Если интеграл $f$ определён, то $\int_E\alpha f\,d\mu=\alpha\int_E f\,d\mu$, $\alpha\in\mathbb R$, с соглашением $0\cdot(\pm\infty)=0$.
  \item Если интеграл $f$ определён, то $|\int_E f\,d\mu|\le\int_E|f|\,d\mu$.
  \item Если $\mu(E)<\infty$ и $a\le f\le b$ на $E$, $a,b\in\mathbb R$, то $f$ суммируема и
  \[
    a\mu(E)\le\int_E f\,d\mu\le b\mu(E).
  \]
  \item Суммируемая на $E$ функция конечна почти везде на $E$.
\end{enumerate}
\end{theorem}

\begin{proof}
\textbf{1.} Для $0\le f\le g$ всякая ступенчатая $s$ с $0\le s\le f$ также не превосходит $g$.
Переход к супремумам даёт неравенство интегралов.
Для суммируемых функций используем $f^+\le g^+$ и $f^-\ge g^-$ и вычитаем конечные интегралы.

\textbf{2.} Равенства следуют из определения интеграла ступенчатой функции.

\textbf{3.} При $f\ge0$ интеграл каждой ступенчатой $s\le f$ по $E$ равен нулю, так как все части разбиения имеют меру нуль.
Поэтому их супремум равен нулю.
Для произвольного $f$ применяем это к $f^+$ и $f^-$.

\textbf{4.} При $\alpha>0$ соответствие $s\mapsto\alpha s$ переводит ступенчатые приближения $f\ge0$ во все ступенчатые приближения $\alpha f$; их интегралы умножаются на $\alpha$.
Для функций произвольного знака применяем формулу к срезкам.
При $\alpha<0$ дополнительно используем $(-f)^+=f^-$, $(-f)^-=f^+$; случай $\alpha=0$ следует из пункта 2.

\textbf{5.} Если $\int_E|f|\,d\mu=\infty$, оценка очевидна.
Иначе $f^\pm\le|f|$ имеют конечные интегралы, и по монотонности из $-|f|\le f\le|f|$ следует требуемое.

\textbf{6.} Срезки $f^\pm$ ограничены неотрицательными постоянными и потому имеют конечные интегралы.
Затем применяем монотонность к $a\le f\le b$ и пункты 2, 4.

\textbf{7.} Для $f\ge0$ и $E_\infty=\{x\in E:f(x)=+\infty\}$ имеем $f\ge n\chi_{E_\infty}$ при любом $n$.
Значит,
\[
  \int_E f\,d\mu\ge n\mu(E_\infty).
\]
Конечность левой части влечёт $\mu(E_\infty)=0$.
Для суммируемой $f$ произвольного знака применяем это к $f^+$ и $f^-$.
\end{proof}
